% Imporditud: tools/import_eksperimendid.py
% Allikas: Eesti füüsikaolümpiaadi eksperimendi ülesannete
%   kogu (Rael Kalda, 2023), ülesanne Ü137, lahendus L137.
\setAuthor{EFO žürii}
\setRound{lõppvoor}
\setYear{2009}
\setNumber{P E2}
\setDifficulty{4}
\setTopic{elektriahelad}
\setVahendid{taskulambi pirn nominaalpingega $U_0$ = 3,6 V ja nominaalvõimsusega $P_0$ = 1,8 W, patarei (pinge mitte üle 3,3 V), tuntud takisti (R = 1 k$\Omega$), voltmeeter}
\setPunktid{12}
\setAllikas{Eksperimendi kogu Ü137, lahendus lk 104}
\setLahendusseis{toores}

\prob{Pirn}
Materjalide eritakistused sõltuvad temperatuurist ning suurte temperatuurierinevuste korral võivad ka takistuste erinevused olla suured. Mitu korda erineb külma lambi (hõõgniit on toatemperatuuril) takistus nominaalrežiimis põleva lambi nominaaltakistusest?

\hint

\solu
Nominaaltakistuse lihtsalt arvutame: $R_1$ = U02/$P_0$. Kuivõrd see on umbes 3 suurusjärku väiksem, kui takistus R, siis takistiga järjestikku patarei klemmidele ühendatuna on lambil eralduv võimsus nominaalrežiimiga võrreldes tühine ning temperatuur faktiliselt toatemperatuur. Mõõdame sellise ühenduse juures pinged lambil $U_1$ ja takistil $U_2$. Lambi takistus $R_0$ = R $\cdot$ $U_1$/$U_2$. Otsitav suhe k = $R_1$/$R_0$.
\probend
